\documentclass[conference]{IEEEtran}
\IEEEoverridecommandlockouts

\usepackage{cite}
\usepackage{amsmath,amssymb,amsfonts}
\usepackage{algorithm}
\usepackage{algorithmic}
\usepackage{graphicx}
\usepackage{pifont}
\usepackage{makecell}
\usepackage{textcomp}
\usepackage{xcolor}
\usepackage{booktabs}
\usepackage{multirow}
\usepackage{array}
\usepackage{hyperref}
\usepackage{amsthm}
\usepackage{float}
\usepackage{tikz}
\usetikzlibrary{positioning,arrows.meta,shapes.geometric,fit,backgrounds,calc}
\theoremstyle{definition}
\newtheorem{proposition}{Proposition}
\newcommand{\cmark}{\ding{51}}
\newcommand{\xmark}{\ding{55}}

% Override IEEEtran's default ALL-CAPS table captions; render in sentence case per IEEE style.
\usepackage[compatibility=false]{caption}
\captionsetup[table]{singlelinecheck=false, font=footnotesize, labelfont={sc}, textfont=normalfont}

\def\BibTeX{{\rm B\kern-.05em{\sc i\kern-.05em b}\kern-.05em\TeX}}

\begin{document}

\title{Intent-Adaptive Hybrid Retrieval and Heterogeneous-Space Person Fusion for Conversational Video Surveillance}

\author{\IEEEauthorblockN{Sujal M H\IEEEauthorrefmark{1},
Sujnan Kumar\IEEEauthorrefmark{1},
Yashas Shetty\IEEEauthorrefmark{1},
Suhan D Shet\IEEEauthorrefmark{1},
and Ravinarayana B\IEEEauthorrefmark{2}}
\IEEEauthorblockA{\IEEEauthorrefmark{1}Department of Computer Science \& Engineering, \\
Mangalore Institute of Technology \& Engineering, Moodabidre, India \\
\IEEEauthorrefmark{2}Head, Department of Computer Science \& Engineering, \\
Mangalore Institute of Technology \& Engineering, Moodabidre, India \\
Email: \{sujalmh9, sujnankumar439, yashass825, suhanshet2004\}@gmail.com, hodcse@mite.ac.in}}

\maketitle

\begin{abstract}
Hybrid retrieval over surveillance video typically combines a structured metadata filter with a dense semantic index using a \emph{fixed} mixing weight, which catastrophically mis-ranks counting and action queries. We make three novel contributions that, in combination, raise mean Precision@5 (macro-averaged over five query categories) on a 150-query natural-language benchmark from 0.67 (best fixed-$\alpha$ baseline, $\alpha{=}0.3$) to 0.83. \textbf{(C1)} An \emph{intent-adaptive blending function} $\alpha(F,\iota,n_s,n_v)$ that produces a closed-form, per-query mixing weight from the parsed filter $F$, the predicted query intent $\iota$, and the structured/semantic result-set imbalance, with two hard overrides for counting and action queries (Eq.~\eqref{eq:adaptive_alpha}); we show (under empirical concavity assumptions) why fixed-$\alpha$ schemes cannot match the per-category optima simultaneously. \textbf{(C2)} A \emph{heterogeneous-space person descriptor} that fuses a SigLIP visual embedding ($\mathbb{R}^{768}$) and a Sentence-Transformers attribute-text embedding ($\mathbb{R}^{384}$) into a single L2-normalized $\mathbb{R}^{1152}$ vector by weighted concatenation rather than the naive same-space addition used by prior multimodal Re-ID work; an ablation shows a 10.7-pp F$_1$ gain over the best unimodal index and a 2.2-pp gain over a same-space late-fusion variant. \textbf{(C3)} A \emph{conversational alert grammar} that reuses the same structured-output LLM parser to materialize fully-typed alert rules (event, object, zone, time, severity, cooldown, side-effects) from free-text utterances, eliminating manual rule authoring; we report 91\% rule-correctness on a held-out alert-authoring benchmark. The contributions are deployed inside a nine-model real-time pipeline that sustains 2~FPS per camera over a 4-camera deployment with a 109~ms median FAISS lookup over 312k indexed person descriptors.
\end{abstract}

\begin{IEEEkeywords}
Hybrid Retrieval, Intent-Adaptive Ranking, Multimodal Person Re-ID, Heterogeneous Embedding Fusion, Conversational Alert Authoring, Vision--Language Surveillance, Natural Language Querying.
\end{IEEEkeywords}

\section{Introduction}

\IEEEPARstart{N}{atural}-language interfaces over surveillance video are now technically feasible: instance segmentation, person Re-ID, and contrastive vision--language models are mature, and structured-output LLMs can convert a user utterance into a typed query schema. The unsolved problem is \emph{how to rank} the resulting candidate frames: a structured filter (\textit{``persons in Zone~A between 22:00 and 23:00''}) and a semantic vector lookup (\textit{``red jacket loitering near the entrance''}) produce two incommensurable score scales. The dominant fusion strategies in the literature are a fixed convex sum $\alpha S_{\text{struct}} + (1-\alpha) S_{\text{sem}}$ tuned once per system~\cite{b7,b11,b12,b29}, parameter-free Reciprocal Rank Fusion~\cite{b27}, or sequential filter-then-rank pipelines~\cite{b30}. We show empirically (Table~\ref{tab:alpha_baseline}) that no single $\alpha$ simultaneously serves counting, finding, tracking, and action queries: a fixed $\alpha=0.5$ collapses to 0.18~P@5 on action queries and a fixed $\alpha=1.0$ collapses to 0.04~P@5 on find queries, while no fixed value exceeds 0.67 macro-mean P@5 across all five categories.

Two further open problems compound this. The per-person descriptor in prior multimodal Re-ID work~\cite{b22b,b32,b33,b34,b35} either fuses same-space embeddings by addition or trains a joint dual encoder, both requiring dimension-matched backbones or joint contrastive data that surveillance deployments rarely have. And alert authoring in commercial engines (Verkada Command, Avigilon Unity, BriefCam Investigator) requires an IT operator to translate a behavioral description into a typed rule, capping the number of useful rules per deployment.

\subsection*{Contributions}
\begin{enumerate}
\item \textbf{Intent-adaptive hybrid retrieval} (\S\ref{sec:adaptive_alpha}). A closed-form, training-free mixing function $\alpha(F,\iota,n_s,n_v)$ with two hard overrides; an override-gap argument (Proposition~\ref{lem:mono}) under empirical concavity; +0.16 macro-mean P@5 over the best fixed-$\alpha$ baseline.
\item \textbf{Heterogeneous-space person fusion} (\S\ref{sec:fusion}). A weighted-concatenation descriptor that fuses backbones of incompatible dimensionality without joint training; +10.7-pp F$_1$@10 over the strongest unimodal index, +2.2-pp over a same-space late-fusion baseline.
\item \textbf{Conversational alert authoring} (\S\ref{sec:alert_grammar}). A typed grammar reusing the query parser; 91\% full-rule correctness (all 9 fields including cooldown and cadence) on a 60-utterance held-out benchmark.
\item \textbf{An end-to-end open implementation} validating all three contributions on a 4-camera live deployment (\S\ref{sec:eval}).\footnote{Code: \url{https://github.com/sujalmh/intent-adaptive-hybrid-surveillance} (to be released upon acceptance).}
\end{enumerate}

\section{Related Work and Position}

We compare directly with the closest published work along the three contribution axes (Table~\ref{tab:related}).

\textbf{Hybrid structured/semantic retrieval.} Combining dense and structured retrieval is established practice in IR~\cite{b25,b27,b29}, dominantly via fixed-$\alpha$ convex combination~\cite{b29} or parameter-free RRF~\cite{b27}. Query-dependent weighting in IR (learning-to-rank, query-dependent fusion) typically learns weights from large-scale click data on text-only sparse/dense pairs and lacks hard overrides for semantically infeasible branches. Recent analogues HyST~\cite{b30} (sequential LLM-filter then dense) and Agentic-RAG~\cite{b3} (router-based) still do not produce a per-query blended score; surveillance forensic-search systems (VIDEX~\cite{b7}, three-stage DL~\cite{b12}, object-IR~\cite{b11}) and CLIP+RAG pipelines~\cite{b1,b19} use a single retrieval branch. We are not aware of prior surveillance-video work that computes the mixing weight per-query from parsed intent and result-set evidence with category hard overrides.

\textbf{Multimodal person Re-ID.} Heterogeneous-dimensional concatenation of frozen encoders is a known retrieval pattern (e.g., RAG~\cite{b4}); we do not claim it is new in IR. We claim its use as the \emph{primary} surveillance-Re-ID descriptor: prior NL-based person-search~\cite{b32,b33,b34,b35} couples a dual encoder by joint training or by descriptive-fusion over same-dimensional modalities, while same-space late-fusion~\cite{b27} requires dimension-matched encoders or two indices. We are not aware of a surveillance Re-ID system using weighted concatenation of heterogeneous frozen encoders as a single FAISS index calibrated on as few as 50 labeled pairs.

\textbf{Conversational alert authoring.} Existing surveillance products~\cite{b7,b8,b12} expose form-driven rule editors and structured-output LLMs~\cite{b20} are standard for slot-filling; the novelty is the specific reuse of the same query-parsing schema to materialize a 9-tuple typed alert rule end-to-end, rather than a new LLM capability.

\begin{table}[!t]
\caption{Position relative to closest related work along the three novel axes; entries indicate the kind of mechanism used.}
\label{tab:related}
\centering
\scriptsize
\setlength{\tabcolsep}{3pt}
\begin{tabular}{@{}p{2.6cm}p{1.25cm}p{1.45cm}p{1.05cm}@{}}
\toprule
System / Work & \makecell[l]{Per-query\\ adapt.\ $\alpha$} & \makecell[l]{Heterog.-sp.\\ Re-ID fusion} & \makecell[l]{Conv.\ alert\\ authoring} \\
\midrule
Fixed-$\alpha$ hybrid~\cite{b29}     & fixed       & ---            & --- \\
RRF~\cite{b27}                        & rank-only   & rank-only      & --- \\
HyST~\cite{b30}                       & sequential  & ---            & --- \\
Agentic RAG~\cite{b3}                 & route-only  & ---            & --- \\
VIDEX~\cite{b7}                       & \xmark      & \xmark         & \xmark \\
3-Stage DL~\cite{b12}                 & \xmark      & \xmark         & \xmark \\
Object-IR~\cite{b11}                  & \xmark      & \xmark         & \xmark \\
LLM-Surv.~\cite{b1}                   & \xmark      & \xmark         & \xmark \\
CUHK-PEDES~\cite{b32}                 & ---         & joint-trained  & --- \\
CLIP-ReID~\cite{b33}                  & ---         & joint-trained  & --- \\
Img-Txt Fusion~\cite{b34}             & ---         & same-space     & --- \\
TriReID~\cite{b35}                    & ---         & same-space     & --- \\
OSNet~\cite{b22b}                     & ---         & visual-only    & --- \\
SigLIP/CLIP~\cite{b24}                & ---         & same-space     & --- \\
\textbf{This work}                    & \cmark      & \cmark         & \cmark \\
\bottomrule
\end{tabular}
\end{table}

\section{Background and System Substrate}

The novel contributions are deployed inside a nine-model real-time pipeline (Table~\ref{tab:models}). Live RTSP / MJPEG streams (1080p, 2~FPS) and uploaded files are processed by YOLOv11m-seg~\cite{b14}, BoT-SORT~\cite{b22} with OSNet $x1\_0$~\cite{b22b} for MOT/Re-ID, OpenVINO PAR-crossroad-0238~\cite{b15} for seven multi-label attributes, and a CIE~LAB histogram with CIEDE2000~\cite{b16} naming over a 25-class CSS3 palette. Two FAISS \texttt{IndexFlatIP} indices~\cite{b25} hold a scene embedding ($\mathbb{R}^{512}$, OpenCLIP ViT-B/32~\cite{b18b}) and our novel person descriptor ($\mathbb{R}^{1152}$, \S\ref{sec:fusion}). MongoDB stores 13 compound-indexed metadata collections; a FastAPI backend exposes 30+ REST endpoints to a Next.js dashboard. Captions come from a tiered VLM pipeline (Qwen2-VL-2B~\cite{b17} with ViT-GPT2 / zero-shot CLIP fallbacks~\cite{b18,b18b,b19}); the NL parser is a structured-output LLM (GPT-4o-mini~\cite{b20} by default; configurable to local Ollama).

\begin{table}[!t]
\caption{Off-the-shelf components and their roles. None are claimed as contributions.}
\label{tab:models}
\centering
\footnotesize
\begin{tabular}{@{}lll@{}}
\toprule
Component & Role & Output \\
\midrule
YOLOv11m-seg~\cite{b14}     & Detection + masks       & per-frame \\
BoT-SORT~\cite{b22}          & MOT association         & track IDs \\
OSNet $x1\_0$~\cite{b22b}    & Re-ID appearance        & 512-d \\
PAR-crossroad-0238~\cite{b15}& Person attributes       & 7-d (sigmoid) \\
OpenCLIP ViT-B/32~\cite{b18b}& Scene embedding         & 512-d \\
SigLIP-base-p16~\cite{b24}   & Person visual           & 768-d \\
MiniLM-L12-v2~\cite{b23}     & Attribute text          & 384-d \\
Qwen2-VL-2B~\cite{b17}       & Captioning              & text \\
GPT-4o-mini~\cite{b20}       & NL parsing              & JSON \\
\bottomrule
\end{tabular}
\end{table}

\section{C1: Intent-Adaptive Hybrid Retrieval}\label{sec:adaptive_alpha}

\subsection{Problem}
A natural-language query $q$ is parsed into a structured filter $F$ (a Pydantic \texttt{NLFilter} of 16 typed fields) and a semantic-text payload $q_{\text{sem}}$. The structured branch retrieves $R_s = \{r\}$ with scores $S_{\text{struct}}(r)\in[0,1]$ from MongoDB; the semantic branch retrieves $R_v$ with scores $S_{\text{sem}}(r)\in[-1,1]$ from FAISS. The hybrid score is the convex combination
\begin{equation}
S_{\text{hyb}}(r) = \alpha\, S_{\text{struct}}(r) + (1-\alpha)\, S_{\text{sem}}(r),\quad \alpha\in[0,1].
\label{eq:hybrid_score}
\end{equation}
Existing systems pick a single $\alpha$ per deployment. We show this is fundamentally insufficient.

\subsection{Failure of fixed-\texorpdfstring{$\alpha$}{alpha} schemes}
We sweep $\alpha\in\{0.0, 0.3, 0.5, 0.7, 1.0\}$ on the 150-query benchmark (Table~\ref{tab:alpha_baseline}). The collapse is structural, not a tuning artifact. \emph{Counting} queries (\textit{``how many cars in Zone~A''}) demand an exact integer that only the structured branch produces; the semantic score is a continuous similarity in $[-1,1]$ that cannot represent ``equals~5,'' so any weight on the semantic branch dilutes a correct count, and P@5 falls from $0.94$ at $\alpha{=}1.0$ to $0.31$ at $\alpha{=}0.0$. \emph{Action} queries (\textit{``person running''}) have no corresponding field in the 16-field \texttt{NLFilter} schema, so the structured branch returns either an unfiltered scan or empty; weight on the structured branch therefore injects noise, and P@5 falls from $0.69$ at $\alpha{=}0.0$ to $0.00$ at $\alpha{=}1.0$. The two requirements push the optimum to opposite endpoints of $[0,1]$; the best fixed macro-mean is only $0.666$ at $\alpha{=}0.3$, motivating a per-query mixing weight rather than a per-deployment one.

\begin{table}[!t]
\caption{Fixed-$\alpha$ baselines on the 150-query NL benchmark (30 queries per category). \textbf{Mean} is the macro-average across the five categories. No single value dominates; counting and action queries pull in opposite directions (boldface = best per row, italics = worst).}
\label{tab:alpha_baseline}
\centering
\footnotesize
\begin{tabular}{@{}lcccccc@{}}
\toprule
$\alpha$ & Count & Find & Track & Action & Filter & Mean \\
\midrule
0.0 (semantic only)  & \emph{0.31} & 0.74 & 0.79 & \textbf{0.69} & 0.55 & 0.616 \\
0.3                  & 0.48        & 0.76 & 0.80 & 0.58           & 0.71 & \textbf{0.666} \\
0.5                  & 0.69        & 0.71 & 0.72 & \emph{0.18}    & 0.84 & 0.628 \\
0.7                  & 0.86        & 0.62 & 0.61 & 0.07           & 0.88 & 0.608 \\
1.0 (structured only)& \textbf{0.94}& \emph{0.04}& \emph{0.02}& \emph{0.00} & \textbf{0.92} & \emph{0.384} \\
\midrule
\textbf{Adaptive} (Eq.~\eqref{eq:adaptive_alpha}) & \textbf{0.94} & \textbf{0.78} & \textbf{0.81} & \textbf{0.69} & \textbf{0.92} & \textbf{0.828} \\
\bottomrule
\end{tabular}
\end{table}

\subsection{Adaptive blending function}
Let $C:\mathcal{F}\to\{\textsc{cnt},\textsc{act},\textsc{soft}\}$ be the \emph{category function} on parsed filters, defined as
\begin{equation}
C(F) = \begin{cases}
\textsc{cnt}, & \text{if } F \text{ contains a numeric count constraint},\\
\textsc{act}, & \text{else if } F \text{ contains an action-verb token},\\
\textsc{soft}, & \text{otherwise}.
\end{cases}
\label{eq:category}
\end{equation}
We then define
\begin{equation}
\alpha(F,\iota,n_s,n_v) =
\begin{cases}
1.00, & C(F){=}\textsc{cnt}, \\[2pt]
0.10, & C(F){=}\textsc{act}, \\[2pt]
\mathrm{clip}_{[0.10,0.95]}(\alpha_{\text{soft}}), & \text{else},
\end{cases}
\label{eq:adaptive_alpha}
\end{equation}
\begin{equation}
\alpha_{\text{soft}} = 0.6\,\alpha_\iota + 0.2\,\rho_{\text{cplx}} + 0.2\,\rho_{\text{src}},
\label{eq:alpha_soft}
\end{equation}
where \textsc{cnt} (count constraint) and \textsc{act} (action verb) are the two override categories of $C(F)$. The intent prior is $\alpha_\iota \in \{0.9, 0.5, 0.7\}$ for $\iota\in\{\text{count}, \text{track}, \text{find}\}$, query complexity $\rho_{\text{cplx}}=\min(1,|F|/4)$, and structured-evidence ratio $\rho_{\text{src}}=n_s/(n_s+n_v+\varepsilon)$ where $n_s=|R_s|, n_v=|R_v|$, $\varepsilon=10^{-6}$. The two hard overrides encode domain truths: a count constraint (e.g., \textit{``how many cars''}) is unanswerable from semantic similarity, and an action verb (e.g., \textit{``running''}) is unanswerable from the structured filter, which has no \texttt{action} field schema.

\subsection{Why fixed \texorpdfstring{$\alpha$}{alpha} cannot match}
\begin{proposition}[Override gap, under empirical assumptions]
\label{lem:mono}
Let $C_1,C_2$ be two queries with conflicting overrides ($C(F_1)=\textsc{cnt}$, $C(F_2)=\textsc{act}$), and assume the empirical P@5 curves $a\mapsto\mathrm{P@5}(a, C_i)$ are concave on $[0.1,1]$ with maxima at $a=1$ and $a=0.1$ respectively (verified on Table~\ref{tab:alpha_baseline}). For any fixed $\alpha^\star\in[0.1,1]$, the worst-case suboptimality
\begin{align}
\Delta(\alpha^\star) = \max\bigl[\,&\mathrm{P@5}(1, C_1) - \mathrm{P@5}(\alpha^\star, C_1),\notag\\
&\mathrm{P@5}(0.1, C_2) - \mathrm{P@5}(\alpha^\star, C_2)\,\bigr]
\end{align}
satisfies $\Delta(\alpha^\star)\ge \tfrac12\bigl(P_{1,1}-P_{1,0.1}\bigr)$, where $P_{i,a}$ denotes the empirical P@5 of category $C_i$ (first subscript: category index $i\in\{1,2\}$) at mixing weight $\alpha=a$ (second subscript).
\end{proposition}
\begin{proof}[Argument]
Any $\alpha^\star\in[0.1,1]$ is at distance at least $0.45$ from one of the two override points. Combined with the stated concavity, this distance lower-bounds the loss against the maximizer of one of the two categories. The argument is empirical, not a worst-case theoretical guarantee; if the assumed concavity is violated the bound need not hold.
\end{proof}
On our benchmark $P_{1,1}-P_{1,0.1}=0.94-0.31=0.63$, so any fixed $\alpha$ loses $\ge 0.31$ P@5 on at least one category; this matches the empirical 0.94 vs. 0.18 gap in Table~\ref{tab:alpha_baseline}.

\subsection{Soft-zone behavior}
For non-override queries, $\alpha_{\text{soft}}$ interpolates intent prior, complexity, and per-branch evidence (weights $0.6/0.2/0.2$); the clipping interval $[0.10, 0.95]$ prevents either branch from being silenced. Realized $\alpha$ ranges are reported in Table~\ref{tab:alpha}.

\begin{table}[!t]
\caption{Realized $\alpha$ ranges on the 150-query benchmark.}
\label{tab:alpha}
\centering
\footnotesize
\begin{tabular}{@{}lc l@{}}
\toprule
Query Class & Realized $\alpha$ & Effect \\
\midrule
Count constraint & 1.00 & Pure structured (override) \\
Action verb      & 0.10 & Semantic-dominant (override) \\
Count intent     & 0.74--0.86 & Structured-leaning (soft) \\
Find intent      & 0.58--0.70 & Balanced (soft) \\
Track intent     & 0.40--0.52 & Semantic-leaning (soft) \\
\bottomrule
\end{tabular}
\end{table}

\subsection{Post-ranking}
Top-$k$ candidates are then re-ranked by an exponential recency boost $S' = S_{\text{hyb}}\cdot 2^{-\Delta t/h}$, $h{=}24$~h, and Maximal Marginal Relevance~\cite{b26} with $\lambda{=}0.3$.

\section{C2: Heterogeneous-Space Person Fusion}\label{sec:fusion}

\subsection{Problem}
Multimodal person Re-ID~\cite{b32,b33,b34,b35} typically requires either two same-dimensional encoders fused by addition, or a jointly trained dual encoder. In a deployment we must work with whatever pre-trained encoders are available, here SigLIP-base ($\mathbb{R}^{768}$)~\cite{b24} and MiniLM-L12-v2 ($\mathbb{R}^{384}$)~\cite{b23}, whose dimensions do not match.

\subsection{Construction}
For each person ROI we form
\begin{align}
\mathbf{v} &= \mathrm{SigLIP}(\mathrm{Mask}(\mathrm{ROI})) \in\mathbb{R}^{768}, \notag\\
\mathbf{t} &= \mathrm{MiniLM}\bigl(\mathrm{Text}(\hat{y}, C)\bigr) \in\mathbb{R}^{384},
\end{align}
where $\mathrm{Text}(\hat{y}, C)$ is the natural-language sentence assembled from the seven sigmoid attribute predictions $\hat{y}$ and the named color $C$ (e.g., \textit{``person wearing red coat, carrying bag, long sleeves, male''}). The fused descriptor is
\begin{equation}
\mathbf{e}_{\text{fused}} = \frac{[\,\alpha_v\,\mathbf{v}\;\Vert\;\beta_t\,\mathbf{t}\,]}{\bigl\Vert [\,\alpha_v\,\mathbf{v}\;\Vert\;\beta_t\,\mathbf{t}\,]\bigr\Vert_2}\;\in\mathbb{R}^{1152},
\label{eq:fusion}
\end{equation}
with $(\alpha_v,\beta_t)=(0.6, 0.4)$ obtained by coarse grid search on the unit simplex (step $0.1$, $\alpha_v+\beta_t=1$, nine pairs) maximizing F$_1$@10 on a held-out 50-pair calibration set. The single L2 normalization makes the inner product a valid cosine similarity searchable by FAISS \texttt{IndexFlatIP}.

\subsection{Why concatenation, not addition}
Same-space addition $\mathbf{v}+\mathbf{t}$ requires a projection $\mathbf{P}\in\mathbb{R}^{768\times 384}$, is sensitive to relative norm, and collapses two information sources into one coordinate per dimension. Weighted concatenation preserves both encoders' full rank at the cost of one normalization, requires no joint training, and the only learned parameter is the scalar pair $(\alpha_v,\beta_t)$.

\subsection{Ablation}
\begin{table}[!t]
\caption{Person-index ablation. ``Add (proj.)'' projects $\mathbf{t}$ to $\mathbb{R}^{768}$ via a random Gaussian projection then adds; ``Late (RRF)'' performs Reciprocal Rank Fusion~\cite{b27} on two parallel indices; ``Concat (ours)'' is Eq.~\eqref{eq:fusion}.}
\label{tab:fusion_ablation}
\centering
\footnotesize
\begin{tabular}{@{}lccc@{}}
\toprule
Person index & P@5 & R@10 & F$_1$@10 \\
\midrule
Visual only ($\mathbf{v}$, $\mathbb{R}^{768}$) & 0.712 & 0.689 & 0.700 \\
Text only ($\mathbf{t}$, $\mathbb{R}^{384}$)    & 0.645 & 0.612 & 0.628 \\
Add (proj.\ $\mathbf{v}+\mathbf{P}\mathbf{t}$, $\mathbb{R}^{768}$) & 0.768 & 0.741 & 0.754 \\
Late fusion (RRF, two indices) & 0.801 & 0.770 & 0.785 \\
\textbf{Concat (ours), $(0.6,0.4)$, $\mathbb{R}^{1152}$} & \textbf{0.823} & \textbf{0.791} & \textbf{0.807} \\
\bottomrule
\end{tabular}
\end{table}
The concat construction beats the strongest unimodal baseline by 10.7~pp F$_1$@10 and the strongest fusion baseline (RRF) by 2.2~pp, using one FAISS index instead of two. A coarse $(\alpha_v,\beta_t)$ sweep is monotonic around $(0.6,0.4)$: $(0.5,0.5){=}0.785$, $(0.7,0.3){=}0.799$, $(0.4,0.6){=}0.774$ F$_1$@10.

\subsection{Indexing complexity}
Adding a person is $O(d)=O(1152)$; querying $k$ neighbors over $N$ persons with \texttt{IndexFlatIP} is $O(Nd)$. On our hardware this yields a 109~ms median for $N=312{,}544$ descriptors (Table~\ref{tab:query_performance}); the corresponding scene-index lookup mean is 107~ms.

\section{C3: Conversational Alert Authoring}\label{sec:alert_grammar}

\subsection{Problem}
Commercial alert engines~\cite{b7,b8,b12} require an operator to translate a behavioral intention into a typed rule (event family, object class, zone, threshold, time window, severity, cooldown, side-effects). For non-technical operators this is a productivity bottleneck that limits the number of useful rules in a deployment.

\subsection{Grammar}
We define a typed alert rule as a 9-tuple
\[
\mathcal{R} = \langle\, e,\, \mathcal{O},\, \zeta,\, \tau,\, \mathcal{T},\, s,\, c,\, \Sigma,\, \kappa\,\rangle
\]
with $e\in\{\text{crowd}, \text{loiter}, \text{fight}, \text{zone}, \text{count}, \text{time}\}$, $\mathcal{O}\subseteq\mathcal{C}_{\text{det}}$ a subset of detected classes, $\zeta$ a zone label, $\tau\in\mathbb{R}_{\ge 0}$ a numeric threshold, $\mathcal{T}=[t_1,t_2]$ a time window with cross-midnight support, $s\in\{\text{low}, \text{med}, \text{high}\}$, $c\in\mathbb{N}$ the cooldown in seconds, $\Sigma\subseteq\{\text{store\_clip}, \text{snapshot}, \text{push\_ws}\}$ side-effects, and $\kappa\in\mathbb{R}_{>0}$ an evaluation cadence (Hz). The grammar is materialized as a Pydantic schema; the same structured-output LLM used for queries fills it.

\subsection{Authoring example}
The utterance \textit{``alert me when more than five people gather in Zone~A after 10~PM''} produces
\begin{align*}
\mathcal{R} = \langle\,&\text{crowd},\;\{\text{person}\},\;\text{Zone~A},\;5,\;[22{:}00,06{:}00],\\
&\text{med},\;60,\;\{\text{snapshot}, \text{push\_ws}\},\;\text{1Hz}\,\rangle
\end{align*}
without any manual configuration. Six rule families are evaluated per-frame against the live detection stream (Table~\ref{tab:alerts}); each rule is cached from MongoDB with a 5-second TTL.

\begin{table}[!t]
\caption{Rule families and triggers.}
\label{tab:alerts}
\centering
\footnotesize
\begin{tabular}{@{}p{2.4cm} p{5.2cm}@{}}
\toprule
Rule Family & Trigger Condition \\
\midrule
Crowd density & $\#\{\text{persons}\} \ge \tau_{\text{crowd}}$ within camera/zone \\
Loitering & $\Delta t_{\text{track}} \ge 60$~s for any track \\
Fight detection & $\bar{\mathcal{E}}_W \ge 0.03 \wedge \exists i\!\neq\!j: \Vert c_i-c_j\Vert<24$~px \\
Zone occupancy & $\#\{\text{persons in ROI}\} > $ capacity \\
Count rules & $\#\{c\} \;\square\; n$, $\square\in\{=,\ge,\le,>,<\}$ \\
Time-of-day & $t \in [t_{\text{start}},t_{\text{end}}]$ (cross-midnight) \\
\bottomrule
\end{tabular}
\end{table}

\subsection{Complementary statistical anomaly detector}
To surface unknown behaviors we maintain a rolling 7-day baseline of detection counts $\{x_{c,h}^{(d)}\}_{d=1}^7$ per (camera, hour) bucket and emit an event when
\begin{equation}
Z_{c,h} = \frac{x_{c,h}^* - \mu_{c,h}}{\sigma_{c,h}+\varepsilon} \ge \tau_Z,
\label{eq:zscore}
\end{equation}
with $\tau_Z\in\{1.5, 2.0, 3.0\}$ for unusual-object, off-hours, and crowd-spike events (Table~\ref{tab:anomaly_results}).

\begin{table}[!t]
\caption{Z-score anomaly detector (7-day labeled subset).}
\label{tab:anomaly_results}
\centering
\footnotesize
\begin{tabular}{@{}lcccc@{}}
\toprule
Event Type & Precision & Recall & F$_1$ & FPR/day \\
\midrule
\texttt{crowd\_spike}    & 0.872 & 0.811 & 0.840 & 1.4 \\
\texttt{off\_hours}      & 0.804 & 0.866 & 0.834 & 2.1 \\
\texttt{unusual\_object} & 0.713 & 0.792 & 0.751 & 3.7 \\
\bottomrule
\end{tabular}
\end{table}

\section{Implementation Detail (Brief)}

The complete detection-and-indexing loop is given in Algorithm~\ref{alg:index}; the resolution loop that consumes Eqs.~\eqref{eq:adaptive_alpha} and~\eqref{eq:fusion} is Algorithm~\ref{alg:resolve}. The end-to-end pipeline is organized as five dataflow layers (capture, perception, indexing, reasoning, output) in which the three novel components (C1, C2, C3) are the only non-off-the-shelf blocks; all other modules are listed in Table~\ref{tab:models}.

\begin{algorithm}[H]
\caption{Detection and indexing per frame.}
\label{alg:index}
\begin{algorithmic}[1]
\STATE \textbf{Input}: frame $F_t$, timestamp $t$
\STATE $D_t \leftarrow \mathrm{YOLOv11mSeg}(F_t,\tau_{\det}{=}0.5)$
\FOR{each $(b_i,m_i,c_i,p_i)\in D_t$}
    \STATE $\mathbf{q}\leftarrow \mathrm{LABHistogramMode}(F_t[b_i],m_i)$
    \STATE $C\leftarrow \arg\min_{\mathbf{p}_j}\Delta E_{00}(\mathbf{q},\mathbf{p}_j)$
    \IF{$c_i{=}\texttt{person} \wedge w(b_i)\ge 60$}
        \STATE $\hat{y}\leftarrow \sigma(\mathrm{PAR}(\mathrm{Letterbox}(\mathrm{Mask},80,160)))$
        \STATE $\mathbf{e}\leftarrow$ Eq.~\eqref{eq:fusion}; $\mathrm{FAISS}_p$.add$(\mathbf{e})$
    \ENDIF
    \STATE MongoDB.insert(record)
\ENDFOR
\IF{$t\bmod\Delta_{\text{idx}}=0$} \STATE EnqueueSemanticIndexing($F_t$) \ENDIF
\end{algorithmic}
\end{algorithm}

\begin{algorithm}[H]
\caption{NL query resolution using Eqs.~\eqref{eq:adaptive_alpha},~\eqref{eq:fusion}.}
\label{alg:resolve}
\begin{algorithmic}[1]
\STATE \textbf{Input}: NL query $q$
\STATE $F, q_{\text{sem}}, \iota \leftarrow \mathrm{LLM\_Parse}(q)$
\STATE $R_s\leftarrow$ MongoDB.find$(F)$; $R_v\leftarrow$ FAISS.search$(g_T(q_{\text{sem}}),k)$ \COMMENT{$g_T$ is the OpenCLIP ViT-B/32 text encoder paired with the scene index}
\STATE $\alpha\leftarrow \alpha(F,\iota,|R_s|,|R_v|)$ \COMMENT{Eq.~\eqref{eq:adaptive_alpha}}
\STATE $R\leftarrow \mathrm{Merge}(R_s,R_v)$; rescore by Eq.~\eqref{eq:hybrid_score}
\STATE $R\leftarrow \mathrm{RecencyBoost}\circ\mathrm{MMR}\circ\mathrm{ConfFilter}(R)$
\RETURN $R[:k]$, $\mathrm{LLM\_Answer}(q,R)$
\end{algorithmic}
\end{algorithm}

\section{Evaluation}\label{sec:eval}

\subsection{Setup}
Hardware: Intel i7-13700K, RTX~4070 (12~GB), 32~GB DDR5, Windows~11. Software: Python~3.12, PyTorch~2.3, CUDA~12.1, OpenVINO~2024.2, FAISS~1.8, MongoDB~7. Live evaluation: four 1080p IP cameras at 2~FPS for $\sim$8~h ($\sim$230k frames, 1.25M detections). Retrieval evaluation: a held-out 150-query NL benchmark spanning five categories (30 queries each) with human-annotated relevant clip sets, plus a 60-utterance held-out alert-authoring benchmark with reference rule annotations.

\textbf{Statistical protocol.} We report mean P@5 with a 95\% bootstrap CI (10{,}000 query-level resamples). Paired comparisons (adaptive vs.\ fixed-$\alpha$, concat vs.\ each fusion baseline) use a two-sided paired bootstrap on per-query P@5 deltas plus a Wilcoxon signed-rank cross-check; significance threshold $p<0.05$.

\subsection{C1 -- adaptive-\texorpdfstring{$\alpha$}{alpha} vs.\ fixed-\texorpdfstring{$\alpha$}{alpha}}
Adaptive-$\alpha$ raises macro-mean P@5 from the best fixed-$\alpha$ value $0.666\,(\pm0.039)$ at $\alpha{=}0.3$ to $0.828\,(\pm0.034)$, an absolute gain of $+0.162$ with paired-bootstrap 95\% CI $[+0.118,+0.207]$ over 150 queries (Table~\ref{tab:alpha_baseline_ci}). The paired bootstrap and Wilcoxon signed-rank tests both reject equality at $p<10^{-4}$; per-category gains against the per-category best fixed-$\alpha$ are zero on the override categories (count, action) and on filter (where $\alpha{=}1.0$ already saturates), and positive on find ($+0.02$, $p=0.04$) and track ($+0.01$, $p=0.07$). Crucially, the gain in macro-mean comes from the fact that no single fixed $\alpha$ can match the count, action, and filter optima simultaneously (Proposition~\ref{lem:mono}); adaptive-$\alpha$ does so by construction.

\begin{table}[!t]
\caption{Per-category P@5 with 95\% bootstrap CIs (10{,}000 resamples, $n{=}30$ per category, $n{=}150$ overall). $\Delta$ is the per-query mean of (adaptive $-$ per-category best-fixed-$\alpha$); $p$ is two-sided paired bootstrap. The bottom row compares the best single fixed-$\alpha$ macro-mean ($\alpha{=}0.3$) against adaptive-$\alpha$.}
\label{tab:alpha_baseline_ci}
\centering
\scriptsize
\setlength{\tabcolsep}{3pt}
\resizebox{\columnwidth}{!}{%
\begin{tabular}{@{}lccc@{}}
\toprule
Category & Best fixed-$\alpha$ & Adaptive & $\Delta$ ($p$) \\
\midrule
Count   & $0.94\pm0.05$ ($\alpha{=}1.0$) & $0.94\pm0.05$ & $+0.00$ (n.s.) \\
Find    & $0.76\pm0.07$ ($\alpha{=}0.3$) & $0.78\pm0.07$ & $+0.02$ ($0.04$) \\
Track   & $0.80\pm0.06$ ($\alpha{=}0.3$) & $0.81\pm0.06$ & $+0.01$ ($0.07$) \\
Action  & $0.69\pm0.08$ ($\alpha{=}0.0$) & $0.69\pm0.08$ & $+0.00$ (n.s.) \\
Filter  & $0.92\pm0.05$ ($\alpha{=}1.0$) & $0.92\pm0.05$ & $+0.00$ (n.s.) \\
\midrule
\textbf{\makecell[l]{Macro-mean\\ (best single fixed-$\alpha$)}} & $0.666\pm0.039$ & $\mathbf{0.828\pm0.034}$ & $\mathbf{+0.162}$ ($<10^{-4}$) \\
\bottomrule
\end{tabular}}
\end{table}

\subsection{C2 -- fusion ablation}
The heterogeneous-space concat descriptor reaches $\mathrm{F}_1@10=0.807\,(\pm0.024)$, beating the best unimodal index ($0.700\pm0.031$) by $+10.7$~pp (paired bootstrap 95\% CI $[+7.9,+13.4]$, $p<10^{-4}$) and the best fusion baseline RRF ($0.785\pm0.026$) by $+2.2$~pp ($95\%$ CI $[+0.4,+4.1]$, $p=0.018$, Wilcoxon $p=0.022$). The same-space additive baseline (random Gaussian projection) gives $+5.3$~pp over visual-only ($p=0.003$) but remains $5.3$~pp below the concat construction ($p=0.001$). Significance therefore holds against both unimodal and fusion baselines, with the concat-vs-RRF margin being the smallest and least conservative.

\subsection{C3 — alert authoring}
The 60 held-out utterances exercise all nine rule fields (Table~\ref{tab:alert_authoring}); ``Rule-correctness'' requires all nine to be correct. Closed-vocabulary fields (severity, cadence, cooldown, event family) reach $0.97$--$1.00$. The hardest are zone naming ($0.88$, fuzzy references such as \textit{``the back entrance''} vs.\ canonical \texttt{Zone~B}) and time window ($0.92$, cross-midnight intervals such as \textit{``after 10~PM''} parsed as $[22{:}00,06{:}00]$); threshold ($0.93$) misses on implicit quantifiers (\textit{``a few''}), and object-set expansion ($0.95$) under-expands category nouns (\textit{``vehicle''} $\to$ \{\texttt{car}\}). Aggregated, $0.91$ of utterances yield a fully-correct rule on the first attempt; the remaining $9\%$ require a single zone- or threshold-level correction, never a structural rewrite.
\begin{table}[!t]
\caption{Conversational alert authoring on a held-out 60-utterance set. ``Field-level'' = per-field exact match, ``Rule-correctness'' = all 9 fields correct.}
\label{tab:alert_authoring}
\centering
\scriptsize
\setlength{\tabcolsep}{3pt}
\begin{tabular*}{\columnwidth}{@{\extracolsep{\fill}}lcl@{}}
\toprule
Field & Acc.\ & Hardest mistake \\
\midrule
Event family       & 0.97 & loiter vs.\ zone \\
Object set         & 0.95 & ``vehicle'' $\to$ \{car, truck\} \\
Zone               & 0.88 & fuzzy zone names \\
Threshold $\tau$   & 0.93 & implicit (``a few'') \\
Time window        & 0.92 & cross-midnight \\
Severity           & 1.00 & --- \\
Cooldown           & 0.98 & --- \\
Side-effects       & 0.95 & implied snapshot \\
Cadence            & 1.00 & --- \\
\midrule
\textbf{Rule-correctness} & \textbf{0.91} & --- \\
\bottomrule
\end{tabular*}
\end{table}

\subsection{Supporting numbers}
Substrate metrics are summarized in Table~\ref{tab:substrate}; the Z-score anomaly detector (Eq.~\eqref{eq:zscore}) is in Table~\ref{tab:anomaly_results}.
\begin{table}[!t]
\caption{Substrate metrics (YOLOv11m-seg detection, BoT-SORT+OSNet tracking mean over 4 cameras, PAR macro, color naming on 1{,}568 ROIs, Z-score anomaly).}
\label{tab:substrate}
\centering
\footnotesize
\begin{tabular}{@{}llc@{}}
\toprule
Stage & Metric & Value \\
\midrule
Detection (person)   & Precision / F$_1$ / mAP@0.5 & 0.965 / 0.927 / 0.873 \\
Tracking (mean)      & MOTA / IDF1                  & 0.909 / 0.865 \\
Person attributes    & Macro F$_1$ (7 attrs)        & 0.882 \\
Color naming         & CIEDE2000 / RGB-Euclid       & 0.923 / 0.821 \\
NL parsing           & Field-correctness            & 0.94 \\
\bottomrule
\end{tabular}
\end{table}

\begin{table}[!t]
\caption{End-to-end query latency by category. ``Sem.\ only'' isolates the FAISS lookup; ``Total'' includes LLM parsing, hybrid blending, MMR, and clip stitching. Person index has $N{=}312{,}544$ descriptors; person-only median is 109~ms.}
\label{tab:query_performance}
\centering
\scriptsize
\setlength{\tabcolsep}{3pt}
\begin{tabular*}{\columnwidth}{@{\extracolsep{\fill}}lcccc@{}}
\toprule
Category & Parse (ms) & Sem.\ only (ms) & Total (ms) & P@5 \\
\midrule
Count   & 1{,}768 & \phantom{0}94 & 3{,}441 & 0.94 \\
Track   & 2{,}101 & 109           & 5{,}413 & 0.81 \\
Find    & 6{,}405 & 121           & 11{,}983 & 0.78 \\
Action  & 2{,}423 & 137           & 6{,}131 & 0.69 \\
Filter  & 1{,}590 & \phantom{0}82 & 3{,}108 & 0.92 \\
\midrule
\textbf{Mean} & 2{,}857 & \textbf{107} & 6{,}015 & \textbf{0.83} \\
\bottomrule
\end{tabular*}
\end{table}

The 4-camera deployment sustains 2.0~FPS per stream at 4.2~GB resident memory, 7.8~GB VRAM, and 68\% mean GPU utilization (91\% p95) on the RTX~4070.

\section{Discussion and Conclusion}

\textbf{Generalization.} All three contributions are schema-decoupled: C1 needs only a category function $C(\cdot)$ with finitely many overrides; C2 admits any pair of frozen encoders, with only $(\alpha_v,\beta_t)$ learned on a 50-pair calibration set; C3's grammar requires only a typed-schema and trigger-function extension per new event family.

\textbf{Limitations.} (i)~Proposition~\ref{lem:mono} is empirical: it assumes per-category concavity of P@5 in $\alpha$, verified on our benchmark but not guaranteed in general. (ii)~Concat fusion costs $d_v+d_t$ bytes per descriptor; galleries with $N>10^7$ would motivate IVF-PQ over flat IP. (iii)~Alert-grammar correctness depends on the LLM; we observed a 6-pp drop with local 7B models. (iv)~\emph{Demographic bias and privacy.} PAR-crossroad-0238 was trained on a corpus skewed toward adult, daylight imagery; its binary \textit{gender} attribute is a coarse visual proxy. We mitigate by using gender only as a soft cue, providing a switch to disable it, and applying a per-camera retention cap on raw imagery.

\textbf{Conclusion.} The three contributions raise macro-mean P@5 from 0.67 to 0.83, gain 10.7-/2.2-pp F$_1$@10 over unimodal/RRF baselines, and reach 91\% full-rule alert-authoring correctness on a 4-camera, 312k-descriptor deployment, showing that schema-aware fusion closes the gap between dense retrieval and a usable conversational surveillance interface.

\begin{thebibliography}{00}
\fontsize{7}{8.4}\selectfont
\setlength{\itemsep}{-1pt}
\setlength{\parsep}{0pt}
\bibitem{b1} U.~De Silva \emph{et~al.}, ``Large language models for video surveillance applications,'' 2025, arXiv:2501.02850. [Online]. Available: \url{https://arxiv.org/abs/2501.02850}

\bibitem{b3} A.~Singh, A.~Ehtesham, S.~Kumar, T.~T.~Khoei, and A.~V.~Vasilakos, ``Agentic retrieval-augmented generation: a survey on agentic RAG,'' 2025, arXiv:2501.09136. [Online]. Available: \url{https://arxiv.org/abs/2501.09136}

\bibitem{b4} O.~Ovadia, M.~Brief, M.~Mishaeli, and O.~Elisha, ``Fine-tuning or retrieval? Comparing knowledge injection in LLMs,'' in \emph{Proc.\ EMNLP}, Nov. 2024, pp.~237--250, doi: 10.18653/v1/2024.emnlp-main.15.

\bibitem{b7} J.~Jung, S.~Park, H.~Kim, C.~Lee, and C.~Hong, ``AI-driven video indexing for rapid surveillance footage summarization and review,'' in \emph{Proc.\ IJCAI}, Aug. 2024, pp.~8687--8690, doi: 10.24963/ijcai.2024/1009.

\bibitem{b8} L.~Rajendran and R.~S.~Shankaran, ``Bigdata enabled realtime crowd surveillance using artificial intelligence and deep learning,'' in \emph{Proc.\ IEEE BigComp}, Jan. 2021, pp.~129--132, doi: 10.1109/BigComp51126.2021.00032.

\bibitem{b11} S.~Jagtap and N.~Chopade, ``Object-based image retrieval and detection for surveillance video,'' \emph{Int.\ J.\ Electr.\ Comput.\ Eng.}, vol.~14, no.~4, pp.~4343--4351, Aug. 2024, doi: 10.11591/ijece.v14i4.pp4343-4351.

\bibitem{b12} J.-W.~Lee and H.-S.~Kang, ``Three-stage deep learning framework for video surveillance,'' \emph{Appl.\ Sci.}, vol.~14, no.~1, art.\ no.~408, Jan. 2024, doi: 10.3390/app14010408.

\bibitem{b14} Ultralytics, ``YOLO11,'' GitHub, 2024. [Online]. Available: \url{https://github.com/ultralytics/ultralytics}

\bibitem{b15} Y.~Gorbachev \emph{et~al.}, ``OpenVINO deep learning workbench: comprehensive analysis and tuning of neural networks inference,'' in \emph{Proc.\ IEEE/CVF ICCVW}, Oct. 2019, doi: 10.1109/ICCVW.2019.00104.

\bibitem{b16} G.~Sharma, W.~Wu, and E.~N.~Dalal, ``The CIEDE2000 color-difference formula: implementation notes, supplementary test data, and mathematical observations,'' \emph{Color Res.\ Appl.}, vol.~30, no.~1, pp.~21--30, Feb. 2005, doi: 10.1002/col.20070.

\bibitem{b17} P.~Wang \emph{et~al.}, ``Qwen2-VL: enhancing vision-language model's perception of the world at any resolution,'' 2024, arXiv:2409.12191. [Online]. Available: \url{https://arxiv.org/abs/2409.12191}

\bibitem{b18} A.~Radford \emph{et~al.}, ``Learning transferable visual models from natural language supervision,'' in \emph{Proc.\ ICML}, vol.~139, 2021, pp.~8748--8763.

\bibitem{b18b} M.~Cherti \emph{et~al.}, ``Reproducible scaling laws for contrastive language--image learning,'' in \emph{Proc.\ IEEE/CVF CVPR}, Jun. 2023, pp.~2818--2829, doi: 10.1109/CVPR52729.2023.00276.

\bibitem{b19} S.~Saha, G.~Singh, A.~Sharma, and M.~Shah, ``CLIP for surveillance: a comparative evaluation,'' 2023, arXiv:2304.12474. [Online]. Available: \url{https://arxiv.org/abs/2304.12474}

\bibitem{b20} J.~Achiam \emph{et~al.}, ``GPT-4 technical report,'' 2023, arXiv:2303.08774. [Online]. Available: \url{https://arxiv.org/abs/2303.08774}

\bibitem{b22} N.~Aharon, R.~Orfaig, and B.-Z.~Bobrovsky, ``BoT-SORT: robust associations multi-pedestrian tracking,'' 2022, arXiv:2206.14651. [Online]. Available: \url{https://arxiv.org/abs/2206.14651}

\bibitem{b22b} K.~Zhou, Y.~Yang, A.~Cavallaro, and T.~Xiang, ``Omni-scale feature learning for person re-identification,'' in \emph{Proc.\ IEEE/CVF ICCV}, 2019, pp.~3702--3712, doi: 10.1109/ICCV.2019.00380.

\bibitem{b23} N.~Reimers and I.~Gurevych, ``Sentence-BERT: sentence embeddings using Siamese BERT-networks,'' in \emph{Proc.\ EMNLP-IJCNLP}, Nov. 2019, pp.~3982--3992, doi: 10.18653/v1/D19-1410.

\bibitem{b24} X.~Zhai, B.~Mustafa, A.~Kolesnikov, and L.~Beyer, ``Sigmoid loss for language image pre-training,'' in \emph{Proc.\ IEEE/CVF ICCV}, Oct. 2023, pp.~11941--11952, doi: 10.1109/ICCV51070.2023.01100.

\bibitem{b25} J.~Johnson, M.~Douze, and H.~J\'egou, ``Billion-scale similarity search with GPUs,'' \emph{IEEE Trans.\ Big Data}, vol.~7, no.~3, pp.~535--547, Jul. 2021, doi: 10.1109/TBDATA.2019.2921572.

\bibitem{b26} J.~Carbonell and J.~Goldstein, ``The use of MMR, diversity-based reranking for reordering documents and producing summaries,'' in \emph{Proc.\ ACM SIGIR}, Aug. 1998, pp.~335--336, doi: 10.1145/290941.291025.

\bibitem{b27} G.~V.~Cormack, C.~L.~A.~Clarke, and S.~B\"uttcher, ``Reciprocal rank fusion outperforms Condorcet and individual rank learning methods,'' in \emph{Proc.\ ACM SIGIR}, Jul. 2009, pp.~758--759, doi: 10.1145/1571941.1572114.

\bibitem{b29} P.~Lewis \emph{et~al.}, ``Retrieval-augmented generation for knowledge-intensive NLP tasks,'' in \emph{Proc.\ NeurIPS}, vol.~33, 2020, pp.~9459--9474.

\bibitem{b30} J.~Myung, J.~Park, and J.~Han, ``HyST: LLM-powered hybrid retrieval over semi-structured tabular data,'' 2025, arXiv:2508.18048. [Online]. Available: \url{https://arxiv.org/abs/2508.18048}

\bibitem{b32} S.~Li, T.~Xiao, H.~Li, B.~Zhou, D.~Yue, and X.~Wang, ``Person search with natural language description,'' in \emph{Proc.\ IEEE CVPR}, 2017, pp.~1970--1979, doi: 10.1109/CVPR.2017.551.

\bibitem{b33} S.~Li, L.~Sun, and Q.~Li, ``CLIP-ReID: exploiting vision--language model for image re-identification without concrete text labels,'' in \emph{Proc.\ AAAI}, vol.~37, no.~1, 2023, pp.~1405--1413, doi: 10.1609/aaai.v37i1.25225.

\bibitem{b34} X.~Li, H.~Guo, M.~Zhang, and B.~Fu, ``Image--text person re-identification with transformer-based modal fusion,'' \emph{Electronics}, vol.~14, no.~3, p.~525, Jan. 2025, doi: 10.3390/electronics14030525.

\bibitem{b35} Y.~Zhai, Y.~Zeng, D.~Cao, and S.~Lu, ``TriReID: towards multi-modal person re-identification via descriptive fusion model,'' in \emph{Proc.\ ACM ICMR}, 2022, pp.~63--71, doi: 10.1145/3512527.3531397.

\end{thebibliography}

\end{document}
